\chapter{REQUIREMENTS ANALYSIS} \label{c3}
\justifying

\section{Overview}

This chapter details the functional and non-functional requirements for the Sign Language Interpreter system, based on the problem statement and stakeholder needs.

\vspace{1em}
\section{Problem Background}

Approximately 63 million deaf people in India face significant communication challenges in their daily interactions with the hearing world. Sign Language (ASL - American Sign Language) is their primary mode of communication, but most hearing individuals cannot understand it. This creates a persistent accessibility gap in educational institutions, healthcare facilities, government services, workplace environments, and public spaces.

Current limitations include:
\begin{itemize}
    \item Manual interpretation is expensive and not always available
    \item Existing sign language recognition systems are either too complex or require specialized hardware
    \item Most systems have poor accuracy and cannot be used in real-time
    \item No affordable solution exists for everyday use
\end{itemize}

\vspace{1em}
\section{Functional Requirements}

Functional requirements specify what the system must do:

\subsection{FR1: Hand Detection}
\begin{itemize}
    \item System shall detect hand(s) in real-time from webcam feed
    \item System shall extract 21 hand landmarks (keypoints) per detected hand
    \item System shall handle single hand detection with confidence threshold $\ge$ 0.7
    \item System shall draw skeleton overlay on detected hand
\end{itemize}

\subsection{FR2: Gesture Recognition}
\begin{itemize}
    \item System shall recognize all 26 ASL alphabets (A-Z)
    \item System shall classify gestures with accuracy $\ge$ 93\%
    \item System shall provide confidence scores for predictions
    \item System shall handle similar letters (M, N, T) with proper preprocessing
\end{itemize}

\subsection{FR3: Real-Time Processing}
\begin{itemize}
    \item System shall process video at 30 FPS minimum
    \item System shall maintain latency $<$ 100ms per frame
    \item System shall work on standard laptop CPU without GPU
    \item System shall use frame buffer for smooth processing
\end{itemize}

\subsection{FR4: Text Output}
\begin{itemize}
    \item System shall buffer recognized letters into words
    \item System shall support space gesture to confirm words
    \item System shall support delete/backspace gesture to remove characters
    \item System shall display recognized text on screen
    \item System shall show confidence score in real-time
\end{itemize}

\subsection{FR5: Text-to-Speech}
\begin{itemize}
    \item System shall convert recognized text to speech
    \item System shall support offline TTS (pyttsx3)
    \item System shall speak complete words when space gesture is detected
    \item System shall allow toggle for TTS on/off
\end{itemize}

\subsection{FR6: User Interface}
\begin{itemize}
    \item System shall display live webcam feed with hand skeleton overlay
    \item System shall show current prediction and confidence score
    \item System shall display word buffer in real-time
    \item System shall provide clear, intuitive controls (Clear, Copy, Settings)
    \item System shall provide visual feedback for detected gestures
\end{itemize}

\vspace{1em}
\section{Non-Functional Requirements}

Non-functional requirements specify system qualities:

\subsection{NFR1: Performance}

The system targets 93%+ accuracy on 6,500+ test samples, latency under 100ms per frame, sustained throughput of 30 FPS, memory usage under 500MB, and model file size under 100MB.

\subsection{NFR2: Compatibility}
\begin{itemize}
    \item \textbf{OS}: Linux, Windows, macOS
    \item \textbf{Python}: 3.8+
    \item \textbf{Hardware}: Standard laptop with webcam (no GPU required)
    \item \textbf{Web Browser}: For Streamlit UI (Chrome, Firefox, Edge)
\end{itemize}

\subsection{NFR3: Usability}
\begin{itemize}
    \item \textbf{Learning Curve}: Intuitive for non-technical users
    \item \textbf{Setup Time}: $<$ 5 minutes for installation
    \item \textbf{Accessibility}: Clear visual feedback and audio output
    \item \textbf{Feedback}: Real-time display of prediction status
\end{itemize}

\subsection{NFR4: Reliability}
\begin{itemize}
    \item \textbf{Robustness}: Works across different lighting conditions
    \item \textbf{Stability}: Handles variable hand sizes and positions
    \item \textbf{Error Handling}: Graceful degradation when hand not detected
    \item \textbf{Recovery}: System recovers quickly from detection failures
\end{itemize}

\subsection{NFR5: Maintainability}
\begin{itemize}
    \item \textbf{Code Quality}: Well-documented, modular code
    \item \textbf{Testing}: Unit tests for core components
    \item \textbf{Deployment}: Easy to package and distribute
    \item \textbf{Extensibility}: Can add new gestures without major changes
\end{itemize}

\vspace{1em}
\section{Data Requirements}

\subsection{DR1: Training Data}
\begin{itemize}
    \item \textbf{Size}: 6,500+ labeled samples (250 per letter × 26 letters)
    \item \textbf{Format}: CSV with 42 features (21 landmarks × 2 coordinates) + label
    \item \textbf{Source}: Custom collection + Kaggle ASL Alphabet dataset
    \item \textbf{Quality}: Clear hand gestures, adequate lighting
\end{itemize}

\subsection{DR2: Model Requirements}
\begin{itemize}
    \item \textbf{Type}: Random Forest Classifier (primary)
    \item \textbf{Input}: Normalized hand landmark coordinates (42 features)
    \item \textbf{Output}: Predicted letter (A-Z) with confidence score
    \item \textbf{Format}: Pickle file (.pkl) for easy loading
\end{itemize}

\vspace{1em}
\section{Technical Stack}

\subsection{Core Libraries}
\begin{table}[H]
\centering
\begin{tabular}{|L{2cm}|C{2cm}|L{5cm}|}
\hline
\textbf{Library} & \textbf{Version} & \textbf{Purpose} \\
\hline
Python & 3.8+ & Programming language \\
\hline
MediaPipe & 0.8+ & Hand detection \& landmark extraction \\
\hline
OpenCV & 4.5+ & Video capture \& frame processing \\
\hline
scikit-learn & 1.0+ & Random Forest classifier \\
\hline
NumPy & 1.20+ & Numerical computations \\
\hline
Pandas & 1.3+ & Data handling \\
\hline
Streamlit & 1.10+ & Web UI \\
\hline
pyttsx3 & 2.90+ & Text-to-speech \\
\hline
\end{tabular}
\end{table}

\vspace{1em}
\section{Constraints}

The system operates under practical constraints: it detects only one hand at a time, requires adequate lighting above 200 lux, the hand must be positioned 30-100cm from the camera, works best with non-cluttered backgrounds, each gesture must be held for 15+ frames, and the palm should face the camera.

\vspace{1em}
\section{Summary}

This chapter has outlined all functional and non-functional requirements for the Sign Language Interpreter system. The requirements provide a clear specification for the development phase, ensuring the final system meets all stakeholder needs and quality standards.